




% The efficacy of Large Language Models (LLMs) in high-performance kernel generation relies critically on the availability of domain-specific data. Unlike general software engineering, kernel generation requires models to internalize hardware intrinsics, parallel execution semantics, and memory hierarchy constraints. In this section, we survey the data landscape and organize resources into two categories: (1) \textit{Training Corpora}, covering both structured datasets and raw kernel repositories; (2) \textit{Knowledge Bases}, which we identify as essential for grounding RAG systems. 

% % A comprehensive index is provided in Table~\ref{tab:full_resources}.

% \subsection{Training Corpora}
% We posit that the utility of training data has evolved from unstructured scraping to targeted, structure-aware curation. We summarize training corpora by their readiness for model pipelines, contrasting pre-processed instruction datasets with a stratified ecosystem of raw repositories.

% \paragraph{Structured Datasets.}
% Structured datasets represent the highest-value signal for instruction tuning, as they explicitly pair intent with optimization. Several recent contributions substantially lower the barrier for cross-language supervision. The KernelBook dataset~\cite{kernelbook2025} supports cross-language alignment by pairing high-level PyTorch operator specifications with optimized Triton implementations at scale. KernelBench further releases the dataset kernelbench-samples, which contains kernel code snapshots generated by the evaluated model during the iterative optimization process, as well as the corresponding performance profiling data. For instruction tuning in HPC settings, HPC-Instruct~\cite{hpcinstruct2024} provides curated instruction--response pairs spanning CUDA, OpenMP, and MPI tasks, enabling models to learn parallelization logic beyond surface syntax. In addition, the HPC-focused subset of The Stack v2~\cite{lozhkov2024stackv2} offers a large unsupervised corpus when filtered for CUDA-related and Triton-related content, serving as a foundational baseline for pre-training.

% \paragraph{Source Code Repositories.}
% While structured datasets provide supervision, the vast majority of domain knowledge remains latent in open-source repositories. To capture the full kernel development lifecycle, we stratify these repositories into three levels based on their abstraction and semantic role.

% \begin{itemize}
%     \item \textbf{Operator and Kernel Libraries}. The first layer consists of operator-centric kernel repositories, which encode a diverse collection of hardware-aware optimization patterns and thus provide high-quality signals that can be distilled through data cleaning and normalization. A prominent subset targets dense linear algebra, where libraries such as CUTLASS~\cite{nvidia2017cutlass} offer highly optimized matrix multiplication kernels with carefully engineered tiling, pipelining, and memory layouts. Another major category focuses on attention-related operators, which have become central to LLM workloads. Representative examples include FlashAttention~\cite{dao2022flashattention}, FlagAttention~\cite{flagattention2023}, AoTriton~\cite{amd2024aotriton}, and xFormers~\cite{xFormers2022}, each encoding distinct optimization strategies for memory-efficient attention across hardware platforms. In addition, a growing class of repositories provides LLM-oriented operator implementations that target both training and inference workloads such as Liger-Kernel~\cite{linkedin2024liger} and FlagGems~\cite{flaggems2024}. Beyond full-precision kernels, several repositories concentrate on model quantization and inference acceleration. Libraries such as Bitsandbytes~\cite{dettmers2022llmint8}, GemLite~\cite{dropbox_gemlite_2024}, FlashInfer~\cite{ye2025flashinfer}, FBGEMM~\cite{khudia2021fbgemm}, and NVIDIA Transformer Engine~\cite{transformerengine2024} implement low-precision arithmetic, fused operators, and numerically stable mixed-precision pipelines. Collectively, these operator repositories form a rich corpus of reusable optimization patterns spanning matrix computation, attention mechanisms, and quantized inference.
%     \item \textbf{Framework and System Integration Code}. The second layer consists of kernels embedded within deep learning frameworks and runtime systems, where operator implementations are designed to support end-to-end usability rather than standalone performance. Core frameworks such as TensorFlow~\cite{abadi2016tensorflowlargescalemachinelearning}, PyTorch (ATen)~\cite{paszke2019pytorchimperativestylehighperformance}, and PaddlePaddle~\cite{ma2019paddlepaddle} provide extensive collections of operator implementations that balance generality, extensibility, and performance. Beyond general-purpose frameworks, inference systems such as vLLM~\cite{kwon2023efficient}, SGLang~\cite{zheng2023sglang}, llama.cpp~\cite{llama_cpp}, and TensorRT-LLM~\cite{nvidia2023tensorrtllm} incorporate specialized kernels tightly coupled with system-level optimizations, for example, dynamic batching and KV-cache management. Similarly, training frameworks such as DeepSpeed~\cite{rasley2020deepspeed} introduce custom operators to support large-scale model parallelism and memory-efficient training. After appropriate cleaning, these framework-level kernels expose how operator implementations are adapted to practical runtime constraints and thus serve as valuable training data.
%     \item \textbf{Domain-Specific Languages (DSLs) and Emerging Abstractions}. The third layer comprises domain-specific languages (DSLs) and their associated tutorials and reference implementations, which offer an intermediate abstraction between high-level intent and low-level kernel code. These resources often include pedagogical examples and reusable operator templates that complement large-scale repositories. DSLs such as Triton, along with emerging languages and abstractions like TileLang and cuTile, provide both language specifications and concrete kernel implementations that illustrate canonical optimization patterns. Tutorial code, sample libraries, and reference operators written in these DSLs can be directly harvested as kernel corpora to learn kernel construction principles in a structured manner. 
% \end{itemize}

% \subsection{Knowledge Bases}

% Beyond executable code, curated domain knowledge also plays a critical role in training kernel-aware models. Such knowledge can be distilled into pre-training corpora to enrich model understanding, or integrated as external knowledge bases to support retrieval-augmented generation in agent-based systems. We argue that such knowledge corpora must be grounded in authoritative specifications to minimize hallucinations. For example, the \textit{CUDA C++ Programming Guide}~\cite{nvidia2024cudaguide} and \textit{PTX ISA Reference}~\cite{nvidia2024ptx} define legal syntax, while \textit{NVIDIA Architecture Tuning Guides}~\cite{nvidia2024tuning} summarize optimization heuristics. However, formal docs often omit engineering ``tribal knowledge''. Thus, community indices such as GPU-MODE~\cite{gpumode2025community} and Triton Index~\cite{tritonindex2025} are valuable for debugging, and meta-indices (e.g., Awesome-CUDA) connect isolated topics. Pedagogical resources like LeetCUDA~\cite{leetcuda2024} and Triton Puzzles~\cite{rush2023tritonpuzzles} aid in step-wise reasoning. Colfax Research~\cite{colfax_research} serves as a specialized technical hub dedicated to High-Performance Computing (HPC) and AI, offering in-depth articles and tutorials on NVIDIA GPU architectures, CUDA optimization, and deep learning frameworks. Finally, profiling outputs from Nsight Compute~\cite{nvidia2024nsight} and Triton-Viz~\cite{tritonviz2024} serve as grounded descriptions of runtime behavior, bridging the gap between static code and dynamic performance.



\begin{table*}[h!]
\centering
% \scriptsize % 使用较小字体以容纳所有内容
% \renewcommand{\arraystretch}{1.15} % 稍微增加行间距，提升密集信息的易读性
% 定义列宽：L=左对齐, p{宽度}=自动换行
\begin{tabular}{l l p{6.8cm} l}
\toprule
\textbf{Data} & \textbf{Resource} & \textbf{Description} & \textbf{Access} \\ 

% ---------------------------------------------------------
% SECTION 1: DATASETS
% ---------------------------------------------------------
\midrule
\multicolumn{4}{l}{\textit{\textbf{I. Structured Datasets (Hugging Face \& Benchmarks)}}} \\ 
\midrule
02/2024 & \textbf{The Stack v2}~\cite{lozhkov2024stackv2} & Unsupervised CUDA/Triton Corpus & \href{https://huggingface.co/datasets/bigcode/the-stack-v2}{[Data]} \\
06/2024 & \textbf{HPC-Instruct}~\cite{hpcinstruct2024} & Instructions for CUDA/MPI/OpenMP & \href{https://huggingface.co/datasets/hpcgroup/hpc-instruct}{[Data]} \\
05/2025 & \textbf{KernelBook}~\cite{kernelbook2025} & Torch-Triton Aligned Corpus & \href{https://huggingface.co/datasets/GPUMODE/KernelBook}{[Data]} \\
02/2025 & \textbf{KernelBench samples} & Kernel Code Snapshots and Profiling Data & \href{https://huggingface.co/datasets/ScalingIntelligence/kernelbench-samples}{[Data]} \\

% ---------------------------------------------------------
% SECTION 2: CODE CORPORA
% ---------------------------------------------------------
\midrule
\multicolumn{4}{l}{\textit{\textbf{II. Code-Centric Corpora (GitHub Repositories)}}} \\
\multicolumn{4}{l}{\cellcolor{gray!10}\textit{Layer 1: High-Performance Operator Libraries }} \\ 
12/2017 & \textbf{CUTLASS} & CUDA C++ Template Library for Matrix Ops & \href{https://github.com/NVIDIA/cutlass}{[Code]} \\
05/2022 & \textbf{FlashAttention}~\cite{dao2022flashattention} & Fast and Memory-Efficient Exact Attention & \href{https://github.com/Dao-AILab/flash-attention}{[Code]} \\
11/2023 & \textbf{FlagAttention}~\cite{flagattention2023} & Memory Efficient Attention Operators in Triton & \href{https://github.com/flagos-ai/FlagAttention}{[Code]} \\
02/2024 & \textbf{AoTriton}~\cite{amd2024aotriton} & AOT-compiled Triton kernels for AMD ROCm & \href{https://github.com/ROCm/aotriton}{[Code]} \\
11/2021 & \textbf{xFormers}~\cite{xFormers2022} & Hackable and Optimized Transformer Blocks & \href{https://github.com/facebookresearch/xformers}{[Code]} \\
08/2024 & \textbf{Liger-Kernel}~\cite{linkedin2024liger} & Efficient Triton Kernels for LLM Training & \href{https://github.com/linkedin/Liger-Kernel}{[Code]} \\
04/2024 & \textbf{FlagGems}~\cite{flaggems2024} & Triton-based Operator Library for LLMs & \href{https://github.com/FlagOpen/FlagGems}{[Code]} \\
09/2022 & \textbf{Bitsandbytes}~\cite{dettmers2022llmint8} & K-bit Quantization Kernels for LLMs & \href{https://github.com/bitsandbytes-foundation/bitsandbytes}{[Code]} \\
09/2024 & \textbf{Gemlite}~\cite{dropbox_gemlite_2024} & Low-Bit Matrix Multiplication Triton Kernels & \href{https://github.com/dropbox/gemlite}{[Code]} \\
01/2025 & \textbf{FlashInfer}~\cite{ye2025flashinfer} & Kernel Library for Efficient LLM Serving & \href{https://github.com/flashinfer-ai/flashinfer}{[Code]} \\
05/2021 & \textbf{FBGEMM}~\cite{khudia2021fbgemm} & Low-Precision Matrix Multiplication & \href{https://github.com/pytorch/FBGEMM}{[Code]} \\
09/2022 & \textbf{Transformer Engine}~\cite{transformerengine2024} & Acceleration Library for Transformer Models & \href{https://github.com/NVIDIA/TransformerEngine}{[Code]} \\

\multicolumn{4}{l}{\cellcolor{gray!10}{\textit{Layer 2: Framework \& System Integration}}} \\
10/2016 & \textbf{PyTorch (ATen)} & Foundational Tensor Library for C++ and Python & \href{https://github.com/pytorch/pytorch}{[Code]} \\
06/2023 & \textbf{vLLM} & High-Efficient Serving Engine & \href{https://github.com/vllm-project/vllm}{[Code]} \\
12/2023 & \textbf{SGLang} & Structured Generation Language for LLMs & \href{https://github.com/sgl-project/sglang}{[Code]} \\
03/2023 & \textbf{llama.cpp} & LLM Inference in C/C++ & \href{https://github.com/ggerganov/llama.cpp}{[Code]} \\
08/2023 & \textbf{TensorRT-LLM} & TensorRT Toolbox for LLM Inference & \href{https://github.com/NVIDIA/TensorRT-LLM}{[Code]} \\
10/2019 & \textbf{DeepSpeed} & System for Large Scale Model Training & \href{https://github.com/deepspeedai/DeepSpeed}{[Code]} \\

\multicolumn{4}{l}{\cellcolor{gray!10}{\textit{Layer 3: Domain-Specific Languages}}} \\
07/2019 & \textbf{Triton} & Open-Source GPU Programming Language & \href{https://github.com/triton-lang/triton}{[Code]} \\
04/2024 & \textbf{TileLang} & Tile-based Optimization Language & \href{https://github.com/tile-ai/tilelang}{[Code]} \\
12/2025 & \textbf{cuTile} & NVIDIA's DSL for Tile-centric Programming & \href{https://docs.nvidia.com/cuda/cutile-python/}{[Link]} \\

% ---------------------------------------------------------
% SECTION 3: EDUCATIONAL
% ---------------------------------------------------------
\midrule
\multicolumn{4}{l}{\textit{\textbf{III. Knowledge Bases \& Educational Resources}}} \\ 
\multicolumn{4}{l}{\textit{Documentation \& Guides}} \\
06/2007 & \textbf{CUDA Guide} & CUDA C++ Programming Guide & \href{https://docs.nvidia.com/cuda/cuda-c-programming-guide/}{[Docs]} \\
06/2007 & \textbf{PTX ISA} & PTX ISA Reference & \href{https://docs.nvidia.com/cuda/parallel-thread-execution/}{[Docs]} \\
05/2020 & \textbf{Tuning Guides} & NVIDIA Architecture Tuning Guides & \href{https://docs.nvidia.com/cuda/}{[Docs]} \\

\multicolumn{4}{l}{\textit{Community Indices \& Tutorials}} \\
01/2024 & \textbf{GPU-MODE} & Resource Stream \& KernelBook & \href{https://github.com/gpu-mode/resource-stream}{[List]} \\
01/2024 & \textbf{Triton Index} & Community Index for Triton Optimization & \href{https://github.com/gpu-mode/triton-index}{[List]} \\
06/2016 & \textbf{Awesome-CUDA} & Community Curated List for CUDA & \href{https://github.com/Erkaman/Awesome-CUDA}{[List]} \\
12/2023 & \textbf{Awesome-GPU} & Awesome GPU Engineering List & \href{https://github.com/goabiaryan/awesome-gpu-engineering}{[List]} \\
05/2023 & \textbf{LeetCUDA} & CUDA Programming Exercises & \href{https://github.com/xlite-dev/LeetCUDA}{[Code]} \\
01/2023 & \textbf{Triton-Puzzles} & Puzzles for Learning Triton & \href{https://github.com/srush/Triton-Puzzles}{[Code]} \\
01/2011 & \textbf{Colfax Research} & Technical Hub Dedicated to HPC and AI & \href{https://research.colfax-intl.com/}{[Link]} \\
09/2018 & \textbf{Nsight Compute} & Kernel Profiling Guide & \href{https://docs.nvidia.com/nsight-compute/}{[Docs]} \\

\bottomrule
\end{tabular}
\caption{A structured overview of training corpora and kernel knowledge bases.  Note that the dates in the table correspond to the initial release; the libraries themselves continue to undergo active development. }
\label{tab:full_resources}
\end{table*}

The efficacy of Large Language Models (LLMs) in high-performance kernel generation relies critically on the availability of domain-specific data. Unlike general software engineering, kernel generation requires models to internalize hardware intrinsics, parallel execution semantics, and memory hierarchy constraints. In this section, we survey the data landscape and organize resources into two categories: (1) \textit{Training Corpora}, covering both structured datasets and raw kernel repositories; (2) \textit{Knowledge Bases}, which we identify as essential for grounding RAG systems. 

The training data consists of targeted, structure-aware curation and unstructured repositories. Structured datasets represent the highest-value signal for instruction tuning, as they explicitly pair intent with optimization. Open-source repositories contain the vast majority of domain knowledge, where optimized kernel code can be extracted and cleaned from open-source operator and kernel libraries, integrated framework or system, and domain-specific languages' tutorials and reference implementations. Beyond executable code, domain knowledge base also plays a critical role in LLM-driven kernel generation. Such knowledge can be distilled into pre-training corpora to enrich model understanding, or integrated as external knowledge bases to support agent-based systems, where the corpora is always provided in authoritative documentation and guides, as well as in community indices or tutorials. A comprehensive index is provided in Table~\ref{tab:full_resources}. Note that the dates listed in the table correspond to the initial release, where these libraries are under active development. 